\section{From approximate midpoints to diamonds}\label{sec:diamonds}

\subsection{An infinite family of approximate midpoints}

Support restrictions disappear when infinitely many separated
perturbations are available. This is the only compactness step in the
metric proof. Approximate-midpoint lenses are central to AMUC
\cite[Theorem~2.1]{DKRRZ2016}; the broader midpoint method is attributed
to Enflo's unpublished argument in \cite[Section~4]{BaudierEtAl2017}.
The configuration below is a separated midpoint family, called a spider
in \cite[Definition~3.7; Definition~6.15 and Proposition~6.16]{BassetLancienProchazka2026}.
We derive its precise squared deficit from Lemma~\ref{lem:paired}.

\begin{lemma}\label{lem:lens}
Let $R>0$, $\varepsilon>0$, $x\in X$, and $(y_j)_{j\ge1}\subset X$ satisfy
\[
 \|x\pm y_j\|\le R,\qquad\|y_i-y_j\|\ge\varepsilon>0
 \quad(i\ne j).
\]
Then
\begin{equation}\label{eq:lens}
 R^2-\|x\|^2\ge\varepsilon^2/48.
\end{equation}
\end{lemma}

\begin{proof}
Convexity gives $\|x\|,\|y_j\|\le R$. Given $\eta>0$, choose
finite ancestral $H$ with $\|Q_Hx\|<\eta$. The bounded sequence
$(P_Hy_j)$ is in a finite-dimensional space, so some distinct $i,j$
satisfy $\|P_H(y_i-y_j)\|<2\eta$. Put $z=(y_i-y_j)/2$.
Since
\[
 x+z=\tfrac12((x+y_i)+(x-y_j)),\quad
 x-z=\tfrac12((x-y_i)+(x+y_j)),
\]
we have $\|x\pm z\|\le R$. Also
\[
 \|P_Hz\|<\eta,\quad\|Q_Hz\|\ge\varepsilon/2-\eta,
 \quad\|P_Hx\|\ge\|x\|-\eta.
\]
Both $P_Hx\pm Q_Hz$ have norm at most $R+2\eta$. By Lemma~\ref{lem:paired},
\[
 (R+2\eta)^2\ge\|P_Hx\|^2+\|Q_Hz\|^2/12.
\]
For $0<\eta<\varepsilon/2$, let $\eta\downarrow0$ to obtain \eqref{eq:lens}.
\end{proof}

\subsection{Adding the loss over the diamond levels}

Figure~\ref{fig:diamond} shows the replacement whose loss we now estimate.
\begin{figure}[t]
\centering
\begin{tikzpicture}[scale=0.85,every node/.style={font=\small}]
 \draw[thick] (0,0)--(2,0);
 \fill (0,0) circle (2pt) node[left] {$u$};
 \fill (2,0) circle (2pt) node[right] {$v$};
 \draw[-{Stealth},thick] (2.8,0)--(4,0);
 \coordinate (u) at (5,0); \coordinate (v) at (9,0);
 \foreach \h/\lab in {1.4/1,0.45/2,-1.4/j}{
   \draw[thick] (u)--(7,\h)--(v);
   \fill (7,\h) circle (2pt) node[above] {$w_{\lab}$};
 }
 \node at (7,-0.5) {$\vdots$};
 \fill (u) circle (2pt) node[left] {$u$};
 \fill (v) circle (2pt) node[right] {$v$};
 \node at (1,-1.9) {one old edge};
 \node at (7,-1.9) {countably many two-edge paths};
\end{tikzpicture}
\caption{One diamond replacement, with only three of the countably many
paths shown. All new edges have length one. The old endpoints are at
distance two, as are any two distinct new middle vertices.}
\label{fig:diamond}
\end{figure}


Each parent edge will consume the same positive amount of squared
stretch. Earlier quantitative diamond arguments accumulate convexity
losses across levels \cite[Lemma~2.2]{LeeNaor2004}, or use a
self-improvement induction \cite[Lemma~1 and Proposition~2]{JohnsonSchechtman2009}.
The countably branching AMUC form appears in
\cite[Section~4]{BaudierEtAl2017}. Here Lemma~\ref{lem:lens} supplies
a fixed squared loss, which we can add directly.

\begin{proof}[Proof of the distortion assertion in Theorem~\ref{thm:main}]

Let $f:D_k\to X$ satisfy $s d_k\le\|f(u)-f(v)\|\le Ds d_k$.
Divide by $s$, so $f$ is noncontracting and $D$-Lipschitz. Distances
between old vertices double at each refinement: a new path between
old vertices consists of two-edge passages which either return to
their starting vertex or contract to an old edge, and every old path
can be lifted with twice its length. Therefore the endpoints of an
edge of $D_m$ have distance $2^{k-m}$ in $D_k$. Distinct middle
vertices replacing that edge also have distance $2^{k-m}$ in $D_k$.

Put
\[
 L_m=\sup_{\{u,v\}\in E(D_m)}
       \frac{\|f(v)-f(u)\|}{2^{k-m}},\qquad 0\le m\le k.
\]
These suprema need not be attained; all that is used is
$1\le L_m\le D$. Fix one parent edge at level $m<k$, put
$\ell=2^{k-m-1}$, and let $w_j$ be its middle vertices. Set
\[
 x=\frac{f(v)-f(u)}{2\ell},\qquad
 y_j=\frac{f(w_j)-(f(u)+f(v))/2}{\ell}.
\]
The child edges give $\|x\pm y_j\|\le L_{m+1}$, and the
middle-vertex distances give $\|y_i-y_j\|\ge2$. Applying Lemma~\ref{lem:lens}
and then taking the supremum over the parent edges yields
\[
 L_m^2\le L_{m+1}^2-1/12.
\]
Summing over $m$ gives
\[
 D^2\ge L_k^2\ge L_0^2+k/12\ge1+k/12.
\]
For $k=0$ this is just $D\ge1$. This completes the proof of
Theorem~\ref{thm:main}.
\end{proof}

To verify the renorming assertion in the introduction, suppose
$a\|z\|\le N(z)\le b\|z\|$. An embedding with distortion at
most $D$ for $N$ has distortion at most $bD/a$ for $\|\cdot\|$.
Thus uniform embeddings for an equivalent norm would contradict the
distortion bound just proved.
